\section{Cơ sở lý thuyết}

Hệ thống truyền phát video trong đồ án được xây dựng dựa trên sự kết hợp chặt chẽ giữa kiến trúc mạng máy tính kinh điển và các kỹ thuật xử lý dữ liệu đa phương tiện thời gian thực, tuân thủ theo các yêu cầu của môn học \cite{lehaminh2024socket}. Phần này sẽ trình bày chi tiết về các giao thức, kỹ thuật socket nâng cao và các mẫu thiết kế phần mềm được áp dụng.

\subsection{Mô hình kiến trúc và Giao diện Socket}

Hệ thống vận hành dựa trên mô hình \textbf{Client-Server}, trong đó trách nhiệm được phân chia rõ ràng giữa hai thực thể. \textbf{Server} đóng vai trò là bên cung cấp dịch vụ thụ động (\textit{Passive Open}), chịu trách nhiệm lưu trữ tài nguyên video, quản lý các phiên làm việc và điều phối luồng dữ liệu. Ngược lại, \textbf{Client} đóng vai trò chủ động (\textit{Active Open}), khởi tạo các yêu cầu kết nối và thực hiện tái tạo hình ảnh từ dữ liệu nhận được.

Để hiện thực hóa cơ chế giao tiếp này, dự án sử dụng \textbf{Socket API} như một giao diện lập trình cốt lõi \cite{stevens2004unix}. Socket đóng vai trò là điểm cuối trong liên kết hai chiều giữa hai chương trình chạy trên mạng. Tùy thuộc vào yêu cầu về độ tin cậy hay tốc độ của từng loại dữ liệu, hệ thống sẽ khởi tạo loại socket phù hợp: \textbf{Stream Socket} (sử dụng TCP) cho kênh điều khiển nhằm đảm bảo tính chính xác tuyệt đối của lệnh, hoặc \textbf{Datagram Socket} (sử dụng UDP) cho kênh dữ liệu để tối ưu hóa độ trễ thời gian thực. Bên cạnh đó, các tùy chọn nâng cao như \texttt{SO\_REUSEADDR} cũng được áp dụng để quản lý trạng thái cổng kết nối hiệu quả, tránh các lỗi khóa tài nguyên thường gặp trong lập trình mạng UNIX \cite{stevens2004unix}.

\subsection{Các giao thức tầng giao vận}

Việc lựa chọn giao thức tầng giao vận là quyết định kỹ thuật quan trọng nhất ảnh hưởng đến hiệu năng của ứng dụng streaming. Thay vì chỉ sử dụng một giao thức đơn lẻ, hệ thống áp dụng cơ chế lai (Hybrid) để tận dụng ưu điểm của cả TCP và UDP theo kiến thức nền tảng về mạng máy tính \cite{kurose2021computer}. Sự khác biệt và phạm vi áp dụng của hai giao thức này trong dự án được tóm tắt tại Bảng \ref{tab:tcp_udp}.

\FloatBarrier
\begin{table}[H]
    \centering
    \caption{So sánh và phạm vi áp dụng của TCP và UDP trong hệ thống}
    \label{tab:tcp_udp}
    \begin{tabular}{p{0.15\textwidth}p{0.45\textwidth}p{0.3\textwidth}}
        \toprule
        \textbf{Giao thức} & \textbf{Đặc điểm kỹ thuật} & \textbf{Áp dụng trong dự án} \\
        \midrule
        \textbf{TCP} & Hướng kết nối, đảm bảo tin cậy, có kiểm soát luồng và tắc nghẽn, độ trễ cao hơn do cơ chế truyền lại (retransmission). & \textbf{Kênh RTSP:} Truyền tải các lệnh điều khiển quan trọng như \texttt{SETUP}, \texttt{PLAY}. \\
        \textbf{UDP} & Phi kết nối, không đảm bảo tin cậy, header nhỏ (8 byte), độ trễ thấp, hỗ trợ truyền thời gian thực. & \textbf{Kênh RTP:} Vận chuyển dữ liệu video (MJPEG). Chấp nhận mất gói để giữ độ trễ thấp. \\
        \bottomrule
    \end{tabular}
\end{table}
\FloatBarrier

\subsection{Giao thức điều khiển RTSP}

\textbf{Real-Time Streaming Protocol (RTSP)}, được định nghĩa trong RFC 2326 \cite{rfc2326}, hoạt động như một "bộ điều khiển từ xa" qua mạng cho các máy chủ đa phương tiện. Điểm khác biệt cốt lõi của RTSP so với HTTP là tính chất \textit{lưu trạng thái} (stateful). Server duy trì một phiên làm việc thông qua mã định danh \textbf{Session ID}, cho phép Server theo dõi trạng thái hiện tại của Client để phản hồi hợp lý. 

Quá trình chuyển đổi trạng thái của hệ thống tuân theo quy trình nghiêm ngặt: từ trạng thái \textbf{INIT} (khởi tạo), chuyển sang \textbf{READY} (sẵn sàng) sau khi thiết lập tài nguyên thành công, và cuối cùng là \textbf{PLAYING} (đang phát) khi luồng dữ liệu bắt đầu được truyền tải. Các phương thức điều khiển chính được mô tả trong Bảng \ref{tab:rtsp_methods}.

\FloatBarrier
\begin{table}[H]
    \centering
    \caption{Các phương thức RTSP và chức năng tương ứng}
    \label{tab:rtsp_methods}
    \begin{tabular}{p{0.2\textwidth}p{0.7\textwidth}}
        \toprule
        \textbf{Phương thức} & \textbf{Mô tả chức năng} \\
        \midrule
        \texttt{SETUP} & Khởi tạo phiên làm việc. Client gửi cổng nhận RTP (\textit{client\_port}), Server phản hồi với \textit{Session ID} và cổng truyền của Server. \\
        \texttt{PLAY} & Yêu cầu Server bắt đầu truyền dữ liệu RTP. Server chuyển trạng thái sang PLAYING. \\
        \texttt{PAUSE} & Tạm dừng luồng dữ liệu. Server ngừng gửi gói tin nhưng vẫn giữ nguyên phiên làm việc và trạng thái kết nối. \\
        \texttt{TEARDOWN} & Kết thúc phiên làm việc. Server đóng kết nối, giải phóng bộ đệm và thu hồi Session ID. \\
        \bottomrule
    \end{tabular}
\end{table}
\FloatBarrier

\subsection{Giao thức vận chuyển RTP và Cấu trúc gói tin}

Trong khi RTSP chịu trách nhiệm điều khiển, \textbf{Real-time Transport Protocol (RTP)} theo tiêu chuẩn RFC 3550 \cite{rfc3550} đảm nhiệm vai trò vận chuyển dữ liệu video thực tế. Do đặc thù của giao thức UDP là không đảm bảo thứ tự, RTP bổ sung một phần tiêu đề (Header) dài 12 byte vào trước mỗi gói dữ liệu để cung cấp thông tin tái đồng bộ.

Ba trường thông tin quan trọng nhất trong RTP Header bao gồm: \textbf{Sequence Number} (16 bit) giúp bên nhận phát hiện mất gói và sắp xếp lại các gói tin bị đảo lộn; \textbf{Timestamp} (32 bit) đóng vai trò nhãn thời gian để đồng bộ hóa tốc độ hiển thị khung hình, đảm bảo video không bị tua nhanh hoặc chậm bất thường; và \textbf{Marker Bit} (1 bit) dùng để đánh dấu gói tin cuối cùng của một khung hình, hỗ trợ quá trình tái lắp ráp dữ liệu phân mảnh.

\subsection{Kỹ thuật Video Streaming và Xử lý bộ đệm}

Để đảm bảo khả năng truyền tải video mượt mà trên môi trường mạng biến động, hệ thống áp dụng các kỹ thuật xử lý dữ liệu nâng cao liên quan đến định dạng nén, phân mảnh và cơ chế đệm.

\subsubsection{Định dạng MJPEG và Phân mảnh gói tin (Fragmentation)}
Dự án sử dụng chuẩn nén \textbf{MJPEG (Motion JPEG)}, nơi mỗi khung hình video là một bức ảnh JPEG độc lập sử dụng kỹ thuật nén nội khung (intra-frame compression) \cite{tekalp2015digital}. Phương pháp này tuân thủ RFC 2435 \cite{rfc2435} và mang lại khả năng chịu lỗi cao, vì việc mất một gói tin chỉ ảnh hưởng cục bộ đến khung hình hiện tại mà không gây lỗi dây chuyền cho các khung hình kế tiếp.

Tuy nhiên, kích thước của một khung hình video (đặc biệt ở độ phân giải HD) thường vượt quá giới hạn \textbf{MTU (Maximum Transmission Unit)} của mạng Ethernet (thông thường là 1500 byte). Để giải quyết vấn đề này, dự án áp dụng kỹ thuật \textbf{Phân mảnh cấp ứng dụng} (Application-Level Fragmentation). Một khung hình lớn sẽ được chia nhỏ thành nhiều gói RTP có kích thước tối đa 1400 byte. Hệ thống sử dụng mẫu thiết kế \textbf{Strategy Pattern} \cite{gamma1994design} để linh hoạt chuyển đổi thuật toán xử lý giữa video SD (có thể không cần phân mảnh) và video HD (bắt buộc phân mảnh), đảm bảo tính tương thích và hiệu năng cao nhất.

\subsubsection{Cơ chế Bộ đệm (Buffering) và Kiểm soát Jitter}
Hiện tượng \textbf{Jitter} (sự biến thiên độ trễ gói tin) là nguyên nhân chính gây ra tình trạng giật hình trong streaming \cite{kurose2021computer}. Để khắc phục, Client được thiết kế hoạt động theo mô hình \textbf{Producer-Consumer} với một \textbf{Frame Buffer} trung gian. Luồng nhận mạng đóng vai trò Producer, liên tục tái lắp ráp các gói tin RTP và đẩy khung hình hoàn chỉnh vào hàng đợi. Luồng hiển thị giao diện đóng vai trò Consumer, lấy khung hình ra để hiển thị theo định kỳ.

Đặc biệt, kỹ thuật \textbf{Pre-buffering} (Tiền tải) được áp dụng nhằm yêu cầu Client phải tích lũy đủ một lượng khung hình nhất định (ngưỡng $N$ frames) trước khi bắt đầu phát. Cơ chế này giúp tạo ra một vùng đệm an toàn, triệt tiêu các dao động tức thời của mạng và đảm bảo trải nghiệm xem video liên tục.

\subsection{Các mẫu thiết kế phần mềm (Design Patterns)}

Để đảm bảo mã nguồn có khả năng mở rộng, dễ bảo trì và tuân thủ các nguyên lý lập trình hướng đối tượng, dự án đã áp dụng một số mẫu thiết kế phần mềm kinh điển được đề xuất bởi Gamma và cộng sự \cite{gamma1994design}. Việc áp dụng các mẫu này giúp giải quyết các vấn đề phổ biến trong việc khởi tạo đối tượng, quản lý trạng thái toàn cục và đơn giản hóa các giao diện phức tạp.

\subsubsection{Mẫu thiết kế Builder (Builder Pattern)}
Trong giao thức RTP, việc khởi tạo một gói tin (\texttt{RTPPacket}) đòi hỏi thiết lập rất nhiều tham số chi tiết cho phần tiêu đề (Header) như: phiên bản (Version), chỉ số đếm (Sequence Number), nhãn thời gian (Timestamp), bit đánh dấu (Marker), và loại tải (Payload Type). Việc sử dụng một hàm khởi tạo (Constructor) với quá nhiều tham số sẽ khiến mã nguồn khó đọc và dễ gây lỗi.

Để giải quyết vấn đề này, dự án áp dụng mẫu thiết kế \textbf{Builder Pattern}. Cụ thể, lớp \texttt{RTPPacketBuilder} (nằm trong thư mục \texttt{common/src/rtp}) được xây dựng để tách biệt quá trình khởi tạo phức tạp khỏi biểu diễn của đối tượng. Mẫu thiết kế này cho phép thiết lập từng thuộc tính của gói tin một cách tuần tự và rõ ràng thông qua chuỗi các phương thức (method chaining) trước khi gọi phương thức \texttt{build()} để tạo ra đối tượng gói tin hoàn chỉnh.

\subsubsection{Mẫu thiết kế Singleton (Singleton Pattern)}
Trong một hệ thống phân tán, có những thành phần cần đảm bảo tính duy nhất và có thể truy cập toàn cục từ bất kỳ đâu trong chương trình, điển hình là bộ ghi nhật ký (Logger) và bộ quản lý cấu hình (Configuration Manager).

Dự án áp dụng mẫu thiết kế \textbf{Singleton Pattern} cho các lớp \texttt{Logger} và \texttt{Config} (trong thư mục \texttt{common/src/utils}). Mẫu thiết kế này đảm bảo rằng chỉ có duy nhất một thể hiện (instance) của lớp \texttt{Logger} được tạo ra trong suốt vòng đời ứng dụng, giúp đồng bộ hóa việc ghi log từ nhiều luồng (thread) khác nhau vào một tệp tin hoặc màn hình console duy nhất, tránh tình trạng xung đột tài nguyên. Tương tự, \texttt{Config} đóng vai trò là kho chứa thông tin cấu hình trung tâm (như địa chỉ IP, cổng, kích thước bộ đệm) mà mọi thành phần khác đều có thể truy xuất.

\subsubsection{Mẫu thiết kế Facade (Facade Pattern)}
Lập trình mạng trên hệ điều hành thường yêu cầu tương tác với các API cấp thấp phức tạp của thư viện Berkeley Sockets (như \texttt{sockaddr\_in}, \texttt{bind}, \texttt{listen}, \texttt{setsockopt}). Việc gọi trực tiếp các hàm C này rải rác trong mã nguồn sẽ làm giảm tính trừu tượng và khó chuyển đổi giữa các nền tảng (như Windows và Linux).

Để khắc phục, lớp \texttt{Socket} (trong thư mục \texttt{common/src/network}) được thiết kế như một \textbf{Facade} (mặt tiền). Lớp này cung cấp một giao diện lập trình hướng đối tượng đơn giản và thống nhất (với các phương thức như \texttt{send()}, \texttt{recv()}, \texttt{bind()}) để che giấu sự phức tạp của các lời gọi hệ thống bên dưới. Nhờ đó, các tầng cao hơn của ứng dụng (như RTSP Server hay RTP Sender) chỉ cần làm việc với đối tượng \texttt{Socket} mà không cần quan tâm đến chi tiết cài đặt của hệ điều hành.

\FloatBarrier
\begin{table}[H]
    \centering
    \caption{Tổng hợp các mẫu thiết kế được áp dụng trong cấu trúc mã nguồn}
    \label{tab:design_patterns}
    \begin{tabular}{p{0.2\textwidth}p{0.3\textwidth}p{0.4\textwidth}}
        \toprule
        \textbf{Mẫu thiết kế} & \textbf{Loại (Category)} & \textbf{Áp dụng thực tế trong dự án} \\
        \midrule
        \textbf{Builder} & Creational (Khởi tạo) & \texttt{RTPPacketBuilder}: Xây dựng gói tin RTP phức tạp từng bước một. \\
        \textbf{Singleton} & Creational (Khởi tạo) & \texttt{Logger}, \texttt{Config}: Quản lý tài nguyên duy nhất (log file, cấu hình) toàn cục. \\
        \textbf{Facade} & Structural (Cấu trúc) & \texttt{Socket}: Đóng gói và đơn giản hóa các thao tác với thư viện socket cấp thấp của hệ điều hành. \\
        \textbf{Strategy} & Behavioral (Hành vi) & \texttt{EncodingStrategy}: Linh hoạt thay đổi thuật toán phân mảnh cho video SD và HD (đã đề cập ở phần trước). \\
        \bottomrule
    \end{tabular}
\end{table}
\FloatBarrier

\subsection{Github Repository}
Dự án được lưu trữ trên Github tại địa chỉ \url{https://github.com/khang1108/rocketProgramming}. Nhóm đã có đẩy code lên và cũng như là tạo các bản \code{release} cho các nền tảng Windows, macOS và Linux.

\section{Tổng quan kiến trúc hoạt động}

\subsection{Sơ đồ triển khai hệ thống (Deployment Diagram)}

Hình \ref{fig:deployment} minh họa kiến trúc triển khai của hệ thống trên hai máy vật lý riêng biệt. Server chạy tiến trình \texttt{./server 8554}, lắng nghe các yêu cầu RTSP trên cổng TCP 8554. Client chạy tiến trình \texttt{./client}, kết nối đến Server để điều khiển phát video và nhận luồng dữ liệu RTP qua UDP. Mỗi thành phần trong sơ đồ đại diện cho một module chức năng cụ thể trong mã nguồn, từ việc xử lý giao thức RTSP, đọc và mã hóa video, đóng gói RTP, đến tái lắp ráp và hiển thị khung hình.

\begin{figure}[H]
\centering
\begin{tikzpicture}[
    node distance=0.8cm,
    machine/.style={draw, rectangle, minimum width=6cm, minimum height=10cm, thick},
    component/.style={draw, rectangle, minimum width=4.5cm, minimum height=1cm, fill=blue!10},
    arrow/.style={->, >=stealth, thick},
    label/.style={font=\small}
]

% Machine 1 - Server
\node[machine] (server_machine) at (0,0) {};
\node[above=0.1cm of server_machine.north, font=\bfseries] {Machine 1 (Server)};
\node[below=0.05cm of server_machine.north, font=\small\ttfamily] {./server 8554};

% Server components
\node[component] (rtsp_server) at (0, 3.5) {RTSPServer};
\node[below=0.1cm of rtsp_server.south, font=\scriptsize] {(TCP:8554)};

\node[component] (video_stream) at (0, 1.8) {VideoStream};
\node[below=0.1cm of video_stream.south, font=\scriptsize] {(Read MJPEG)};

\node[component] (encoding_ctx) at (0, 0.1) {EncodingCtx};
\node[below=0.1cm of encoding_ctx.south, font=\scriptsize] {(Strategy)};

\node[component] (rtp_builder) at (0, -1.6) {RTPPacketBuilder};
\node[below=0.1cm of rtp_builder.south, font=\scriptsize] {(Builder)};

\node[component] (udp_socket_server) at (0, -3.3) {UDP Socket};
\node[below=0.1cm of udp_socket_server.south, font=\scriptsize] {(Send packets)};

% Arrows between server components
\draw[arrow] (rtsp_server) -- (video_stream);
\draw[arrow] (video_stream) -- (encoding_ctx);
\draw[arrow] (encoding_ctx) -- (rtp_builder);
\draw[arrow] (rtp_builder) -- (udp_socket_server);

% Machine 2 - Client
\node[machine] (client_machine) at (10,0) {};
\node[above=0.1cm of client_machine.north, font=\bfseries] {Machine 2 (Client)};
\node[below=0.05cm of client_machine.north, font=\small\ttfamily] {./client ...};

% Client components
\node[component] (rtsp_client) at (10, 3.5) {RTSPClient};
\node[below=0.1cm of rtsp_client.south, font=\scriptsize] {(State Pattern)};

\node[component] (frame_buffer) at (10, 1.8) {FrameBuffer};
\node[below=0.1cm of frame_buffer.south, font=\scriptsize] {(Observer)};

\node[component] (frame_display) at (10, 0.1) {FrameDisplay};
\node[below=0.1cm of frame_display.south, font=\scriptsize] {(UI)};

\node[component] (frame_reassembler) at (10, -1.6) {FrameReassembler};

\node[component] (rtp_receiver) at (10, -3.3) {RTPReceiver};
\node[below=0.1cm of rtp_receiver.south, font=\scriptsize] {(UDP:25000)};

% Arrows between client components
\draw[arrow] (rtsp_client) -- (frame_buffer);
\draw[arrow] (frame_buffer) -- (frame_display);
\draw[arrow] (frame_reassembler) -- (frame_display);
\draw[arrow] (rtp_receiver) -- (frame_reassembler);

% Communication between machines
\draw[<->, thick, red] (rtsp_server.east) -- (rtsp_client.west) node[midway, above, font=\small] {RTSP (TCP)};
\draw[->, ultra thick, green!60!black] (udp_socket_server.east) -- (rtp_receiver.west) node[midway, above, font=\small] {RTP (UDP)};

\end{tikzpicture}
\caption{Sơ đồ triển khai hệ thống video streaming phân tán}
\label{fig:deployment}
\end{figure}

\subsection{Định dạng thông điệp RTSP (RTSP Message Format)}

Giao thức RTSP sử dụng cú pháp văn bản dạng ASCII, tương tự HTTP, để dễ dàng đọc và gỡ lỗi. Mỗi thông điệp RTSP bao gồm một dòng bắt đầu (Request Line hoặc Status Line), các dòng tiêu đề (Headers), và một dòng trống (\texttt{\textbackslash r\textbackslash n}) để đánh dấu kết thúc thông điệp.

\subsubsection{Định dạng RTSP Request}
Một yêu cầu RTSP từ Client gửi đến Server có cấu trúc như sau (xem Hình \ref{fig:rtsp_request}):

\begin{figure}[H]
\centering
\begin{tikzpicture}[
    msgbox/.style={draw, rectangle, minimum width=9cm, minimum height=0.8cm, fill=yellow!10, font=\ttfamily\small, anchor=west},
    annotation/.style={font=\small\itshape, red!70!black}
]

% Draw message box
\draw[thick] (0,0) rectangle (10,3.5);

% Content lines
\node[msgbox] (line1) at (0.2, 2.9) {METHOD FILE RTSP/1.0\textbackslash r\textbackslash n};
\node[annotation, right=0.2cm of line1.east, anchor=west] {$\leftarrow$ Request line};

\node[msgbox] (line2) at (0.2, 2.0) {CSeq: <sequence\_number>\textbackslash r\textbackslash n};
\node[annotation, right=0.2cm of line2.east, anchor=west] {$\leftarrow$ Mandatory header};

\node[msgbox] (line3) at (0.2, 1.1) {[Additional headers]\textbackslash r\textbackslash n};
\node[annotation, right=0.2cm of line3.east, anchor=west] {$\leftarrow$ Optional headers};

\node[msgbox] (line4) at (0.2, 0.2) {\textbackslash r\textbackslash n};
\node[annotation, right=0.2cm of line4.east, anchor=west] {$\leftarrow$ Empty line (end)};

\end{tikzpicture}
\caption{Cấu trúc thông điệp RTSP Request}
\label{fig:rtsp_request}
\end{figure}

Trong đó:
\begin{itemize}
    \item \textbf{METHOD}: Phương thức điều khiển (ví dụ: \texttt{SETUP}, \texttt{PLAY}, \texttt{PAUSE}, \texttt{TEARDOWN})
    \item \textbf{FILE}: Đường dẫn đến tài nguyên video trên Server (ví dụ: \texttt{movie.mjpeg})
    \item \textbf{RTSP/1.0}: Phiên bản giao thức
    \item \textbf{CSeq}: Số thứ tự yêu cầu (bắt buộc), dùng để đối chiếu với phản hồi
    \item \textbf{Additional headers}: Các tiêu đề tùy chọn khác như \texttt{Session}, \texttt{Transport}
\end{itemize}

\subsubsection{Định dạng RTSP Response}
Server phản hồi lại Client với thông điệp có cấu trúc như Hình \ref{fig:rtsp_response}:

\begin{figure}[H]
\centering
\begin{tikzpicture}[
    msgbox/.style={draw, rectangle, minimum width=9cm, minimum height=0.8cm, fill=green!10, font=\ttfamily\small, anchor=west},
    annotation/.style={font=\small\itshape, blue!70!black}
]

% Draw message box
\draw[thick] (0,0) rectangle (10,3.5);

% Content lines
\node[msgbox] (line1) at (0.2, 2.9) {RTSP/1.0 <status\_code> <reason>\textbackslash r\textbackslash n};
\node[annotation, right=0.2cm of line1.east, anchor=west] {$\leftarrow$ Status line};

\node[msgbox] (line2) at (0.2, 2.0) {CSeq: <sequence\_number>\textbackslash r\textbackslash n};
\node[annotation, right=0.2cm of line2.east, anchor=west] {$\leftarrow$ Echo CSeq};

\node[msgbox] (line3) at (0.2, 1.1) {[Additional headers]\textbackslash r\textbackslash n};
\node[annotation, right=0.2cm of line3.east, anchor=west] {$\leftarrow$ Headers};

\node[msgbox] (line4) at (0.2, 0.2) {\textbackslash r\textbackslash n};
\node[annotation, right=0.2cm of line4.east, anchor=west] {$\leftarrow$ Empty line (end)};

\end{tikzpicture}
\caption{Cấu trúc thông điệp RTSP Response}
\label{fig:rtsp_response}
\end{figure}

\noindent
Trong đó:
\begin{itemize}
    \item \textbf{status\_code}: Mã trạng thái HTTP-like (ví dụ: \texttt{200} OK, \texttt{404} Not Found)
    \item \textbf{reason}: Chuỗi mô tả trạng thái (ví dụ: \texttt{OK}, \texttt{Not Found})
    \item \textbf{CSeq}: Phản chiếu lại số thứ tự từ yêu cầu để Client đối chiếu
    \item \textbf{Headers}: Các thông tin bổ sung như \texttt{Session}, \texttt{Transport}, \texttt{Content-Length}
\end{itemize}

Việc sử dụng \texttt{CSeq} (Command Sequence) trong cả Request và Response đảm bảo tính toàn vẹn của giao tiếp, đặc biệt trong môi trường mạng bất đồng bộ, giúp Client xác định chính xác phản hồi nào tương ứng với yêu cầu nào khi có nhiều yêu cầu được gửi đồng thời.

\subsection{Cấu trúc thư mục dự án (Project Folder Structure)}

Dự án được tổ chức theo mô hình module hóa rõ ràng, tách biệt giữa Server, Client và Common library. Hình \ref{fig:folder_structure} minh họa cấu trúc thư mục chi tiết của toàn bộ mã nguồn.

\begin{figure}[H]
\centering
\begin{lstlisting}[basicstyle=\ttfamily\small, frame=single, xleftmargin=1cm]
rocketProgramming/
|-- common/
|   |-- include/     (network, rtp, patterns, utils)
|   |-- src/         (Socket, Logger, Config)
|   `-- tests/       (test_rtsp_message, test_timer)
|
|-- server/
|   |-- include/     (network, rtp, video)
|   |-- src/         (RTSPServer, VideoStream, Encoding)
|   |-- config/      (server.conf)
|   `-- videos/      (*.Mjpeg files)
|
|-- client/
|   |-- include/     (ui, network, rtp, buffer)
|   `-- src/         (ClientUI, RTSPClient, RTPReceiver)
|
|-- setup/           (setup.py, setup.spec)
\end{lstlisting}
\caption{Cấu trúc thư mục dự án với phân chia module rõ ràng}
\label{fig:folder_structure}
\end{figure}

\textbf{Giải thích cấu trúc:}

\begin{itemize}
    \item \textbf{common/}: Thư viện dùng chung cho cả Server và Client
    \begin{itemize}
        \item \texttt{include/}: Header files tổ chức theo namespace (network, rtp, patterns, utils)
        \item \texttt{src/}: Implementation của Socket wrapper, RTP packet, Logger, Config manager
        \item \texttt{tests/}: Unit tests cho các thành phần cốt lõi (RTSP message parsing, Timer)
    \end{itemize}
    
    \item \textbf{server/}: Chương trình Server
    \begin{itemize}
        \item \texttt{include/network/}: RTSPServer, ServerWorker (multi-threading)
        \item \texttt{include/rtp/}: RTP sender, packet builder
        \item \texttt{include/video/}: VideoStream (MJPEG extraction), EncodingStrategy
        \item \texttt{src/}: Implementation tương ứng
        \item \texttt{config/}: File cấu hình server (port, video path, log level)
        \item \texttt{videos/}: Chứa các file video MJPEG để streaming
    \end{itemize}
    
    \item \textbf{client/}: Chương trình Client với Qt6 UI
    \begin{itemize}
        \item \texttt{include/ui/}: ClientUI, BufferedSlider (custom Qt widgets)
        \item \texttt{include/network/}: RTSPClient (state machine)
        \item \texttt{include/rtp/}: RTPReceiver, FrameReassembler
        \item \texttt{include/buffer/}: FrameBuffer (thread-safe queue, prebuffering)
        \item \texttt{src/}: Implementation tương ứng
    \end{itemize}
    
    \item \textbf{setup/}: Scripts tự động hóa build và deployment
    \begin{itemize}
        \item \texttt{setup.py}: Python script kiểm tra dependencies và build (cross-platform)
        \item \texttt{setup.spec}: PyInstaller spec để đóng gói
        \item \texttt{setup\_windows.bat}: Batch script cho Windows
    \end{itemize}
\end{itemize}

Cấu trúc này tuân theo nguyên tắc \textbf{separation of concerns}: mỗi module có trách nhiệm rõ ràng, header và source tách biệt, tests độc lập. Điều này giúp dễ dàng maintain, extend và test từng phần của hệ thống.
